\section{series with positive terms}

\index{series!with positive terms}
In the following, we will consider series whose terms are positive real numbers
(or at least non-negative).
When we write $a_n >0$ (or $a_n \ge 0$), it will always be implied that $a_n\in
\RR$ since for non-real complex numbers we have not defined the order relation.

\begin{theorem}[nature of series with positive terms]\label{th:serie_positiva}
\mymark{***}%
\mymargin{nature of series with positive terms}%
\index{nature!of a series with positive terms}%
If $a_n\ge 0$, the series $\sum a_n$ is regular:
either it converges or it diverges to $+\infty$.
\end{theorem}
%
\begin{proof}
\mymark{***}
If $a_n \ge 0$, since $S_{n+1} = S_n + a_n$, it means that the sequence $S_n$ of
partial sums is increasing.
Thus, the limit of $S_n$ exists and cannot be negative.
\end{proof}

\begin{theorem}[comparison test]
\mymark{**}
\mymargin{comparison test}
\index{test!comparison for series}
Let $\sum a_n$ and $\sum b_n$ be series with positive terms that are compared:
$0\le a_n\le b_n$.
Then
\[
  \sum_{k=0}^{+\infty} a_k \le \sum_{k=0}^{+\infty} b_k.
\]
In particular, if $\sum b_n$ converges, then $\sum a_n$ also converges, and if
$\sum a_n$ diverges, then $\sum b_n$ also diverges.

This latter result holds even if $0 \le a_n \ll b_n$.
\end{theorem}
%
\begin{proof}
\mymark{*}
If $S_n$ are the partial sums of $\sum a_n$ and $R_n$ are the partial sums of
$\sum b_n$, we have $S_n \le R_n$, and the result reduces to the comparison
between sequences.

In the case where $a_n \ll b_n$, by definition, we know that $\frac{a_n}{b_n}\to
0$, and therefore, from the definition of limit, we know that there exists $N$
such that for every $n>N$, we have (having chosen $\eps=1$)
\[
  \frac{a_n}{b_n} < 1.
\]
Thus, we obtain $a_n \le b_n$ for all $n$ except at most a finite number.
Knowing that the nature of the series does not change if the series is modified
on a finite number of terms, we reduce to the previous case.
\end{proof}

\begin{example}\label{ex:52573}
\mymark{***}
The series
\begin{equation}\label{eq:296453}
 \sum_{k=1}^{+\infty} \frac{1}{k^2}
\end{equation}
is convergent.
Indeed, observing that for every $n>0$, we have
\[
  \frac{1}{(n+1)^2} \le \frac{1}{n(n+1)}
\]
we can assert that
\[
  \sum_{k=1}^{+\infty} \frac{1}{k^2}
  = 1 + \sum_{k=1}^{+\infty} \frac{1}{(k+1)^2}
  \le 1+ \sum_{k=1}^{+\infty} \frac{1}{k(k+1)}
  = 2
\]
since we have reduced to the telescopic series of Mengoli, which has a sum equal
to $1$.

We thus know that the series~\eqref{eq:296453} is convergent without knowing
exactly what its sum is.
However, we can numerically find approximations of the sum by summing the first
terms and estimating the error using the series of Mengoli, whose sum we know
how to calculate.
Indeed, if $S_N$ is the partial sum of the first $N$ terms and $S = \lim S_N$ is
the sum of the series, since $1/(k+1)^2 \le 1/(k^2+k)$, we have
\[
S_N
\le S
\le S_N + \sum_{k=N}^{+\infty} \frac{1}{k(k+1)}
\le S_N + \frac{1}{N}.
\]
To calculate the first 6 exact decimal places, it will suffice to sum the first
million terms of the series.
This can be done, for example, with the code provided on
page~\pageref{code:series}, obtaining $S\approx 1.644934$.

Using much more advanced tools, Euler \index{Euler}
(Leonard Euler, 1707--1783) managed to express the sum of this series using
fundamental mathematical constants (we will do so with a different method in
exercise~\ref{ex:Basilea}).
\end{example}

\begin{corollary}[asymptotic comparison test]
\mymark{*}
\mymargin{asymptotic comparison test}
\index{test!asymptotic comparison}
If $a_n$ and $b_n$ are sequences with positive terms, asymptotically equivalent
(definition~\ref{def:ordine_infinito}), then the corresponding series $\sum a_n$
and $\sum b_n$ have the same nature.
\end{corollary}
%
\begin{proof}
\mymark{*}
Series with positive terms cannot be indeterminate, so it is sufficient to
verify that if one series converges, the other also converges.
Since $a_n / b_n$ is convergent, this ratio must also be bounded
(theorem~\ref{th:limitatezza_successione_convergente}), so there exists $C\in
\RR$ such that
\[
   a_n \le C \cdot b_n.
\]
If the series $\sum b_n$ converges, then $\sum C \cdot b_n$ also converges, and
by comparison, $\sum a_n$ also converges.

Conversely, by swapping the roles of $a_n$ and $b_n$, we verify that if $a_n$
converges, then $b_n$ also converges.
\end{proof}

\begin{example}
The series
\[
\sum_n \frac{n^2+2n+3}{2n^4-n^3+n+1}
\]
is convergent. Indeed, it can be easily verified that
\[
   \frac{n^2+2n+3}{2n^4-n^3+n+1} \sim \frac{1}{2n^2}.
\]
But we know that the series $\sum 1/n^2$ is convergent, consequently, the series
$\sum 1/(2n^2)$ is also convergent (by the linearity of the sum,
theorem~\ref{th:linearita_somma_serie}), and therefore, by asymptotic
comparison, the given series is also convergent.
\end{example}

In chapter \ref{sec:serie_segno_variabile}, we will see that the asymptotic
convergence test can only be applied to series with positive terms.

\begin{theorem}[root test]
\mymargin{root test}
\index{test!root}
Let $\sum a_n$ be a series with non-negative terms (i.e., $a_n\ge 0$) such that
\mymark{***}
$\sqrt[n]{a_n} \to \ell \in [0,+\infty]$.
If $\ell<1$, then the series converges.
If $\ell>1$, then the series diverges.

More generally, the result is valid with
\[
  \ell = \limsup \sqrt[n]{a_n}
\]
even in the case where the limit of $\sqrt[n]{a_n}$ does not exist.
\end{theorem}
%
\begin{proof}
\mymark{***}
In the case $\ell < 1$, we take $q$ with $\ell < q < 1$ and set $\eps = q-\ell$.
By the definition of limit $\sqrt[n]{a_n}\to \ell$ (but it suffices that
$\limsup \sqrt[n]{a_n}=\ell$), we know there exists $N$ such that for every $n >
N$, we have
\[
  \sqrt[n]{a_n} < \ell + \eps = q
\]
that is
\[
   a_n < q^n.
\]
Knowing that $\sum q^n$ converges, and knowing that the nature of the series
does not change by modifying a finite number of terms, by comparison, we can
conclude that the series $\sum a_n$ also converges.

If $\ell>1$, we have that $\sqrt[n]{a_n}>1$ and therefore $a_n>1$ for infinitely
many values of $n$. The sequence $a_n$ is not infinitesimal, and thus the series
cannot converge.
\end{proof}

\begin{example}
The series
\[
  \sum_k 2^{(\ln k) - k}
\]
is convergent. Indeed, we have
\[
 \sqrt[k]{2^{\ln k - k}}
 = 2^{\frac{\ln k - k}{k}}
 = 2^{\frac{\ln k }k - 1}
 \to 2^{-1}
 = \frac{1}{2}
 < 1.
\]
\end{example}

\begin{theorem}[ratio test]
\mymark{***}
\mymargin{ratio test}
\index{test!ratio for series}
Let $\sum a_n$ be a series with positive terms ($a_n > 0$) such that $a_{n+1} /
a_n \to \ell \in [0,+\infty]$.
If $\ell <1$, then the series converges.
If $\ell > 1$, the series diverges.
\end{theorem}
%
\begin{proof}
\mymark{*}
It would not be difficult to provide a direct proof, similar to the proof given
for the root test.
However, we can observe that by the Cesàro convergence criterion
(theorem~\ref{th:criterio_cesaro}), we have $\sqrt[n]{a_n} \to \ell$, so we
reduce to the root test without needing further proofs.
\end{proof}

\begin{example}
\mymark{***}
For every $x\ge 0$, the series
\[
  \sum \frac{x^n}{n!}
\]
converges.
\end{example}
%
\begin{proof}
We apply the ratio test.
Setting $a_n = x^n / n!$, we have
\[
\frac{a_{n+1}}{a_n}
= \frac{x^{n+1}}{(n+1)!}\cdot \frac{n!}{x^n}
= \frac{x}{n+1} \to 0 < 1.
\]
Thus, the series converges.
\end{proof}

\subsection{associativity of the sum of a series}

\begin{example}
Setting $a_k = (-1)^k$, we know that the corresponding series:
\[
  1 - 1 + 1 - 1 + 1 \dots
\]
is indeterminate because the partial sums are
\[
  S_n = \sum_{k=0}^n (-1)^k
  = \begin{cases}
   1 & \text{if $n$ is even}\\
   0 & \text{if $n$ is odd.}
  \end{cases}
\]
However, if we associate the terms in pairs, we obtain the series
\[
  (1-1) + (1-1) + (1-1) \dots = 0 + 0 + 0 + \dots = 0
\]
which is convergent.
The reason is that the sequence of partial sums of this new series is a
particular subsequence (the one with even indices) of the sequence of partial
sums of the original series. Therefore, it should not surprise us that even
though the original series was indeterminate, it is possible to associate the
terms of the series in such a way as to obtain a regular series (in this case,
convergent).
In fact, by the Bolzano-Weierstrass theorem~\ref{th:bolzano_weierstrass}, we
know that it is always possible to extract a regular (convergent or divergent)
subsequence from any sequence.
Even when a convergent series is extracted from an indeterminate series, the sum
of the series can depend on how the terms are associated.
For example, in the previous series, if we only took the sums with odd indices,
we would obtain:
\[
  1 + (-1+1) + (-1+1) + \dots = 1.
\]
\end{example}

The previous example shows us that by associating the terms of an indeterminate
series, it is possible to obtain series with different nature and sum.
The following theorem tells us that this "bad" phenomenon can only occur when
starting from an indeterminate series.
If the series is regular, then we can associate its terms without changing
either its nature or its sum.
In particular, this is true for series with positive terms.

\begin{theorem}[associativity of regular series]%
\label{th:serie_associativa}%
If $\sum a_k$ is a series, choosing any increasing sequence $k_n$ with $k_0=0$,
we can consider the series $\sum b_n$ whose terms
\[
  b_n = \sum_{j=k_n}^{k_{n+1}-1} a_j
\]
are obtained by associating the terms of $a_k$ into consecutive groups delimited
by the sequence of indices $k_n$.

If the series $\sum a_k$ is regular (convergent or divergent), then the series
$\sum b_n$ is also regular, and we have
\[
\sum_{n=0}^{+\infty} b_n
= \sum_{k=0}^{+\infty} a_k.
\]

In particular, this holds if the series $\sum a_k$ is with positive terms.
\end{theorem}
%
\begin{proof}
Let $S_k = \sum_{j=0}^k a_j$ be the partial sums of the series $\sum a_j$. Then
the partial sums of the series $\sum b_n$ are nothing but the subsequence
$S_{k_n}$.
Thus, if $S_k$ converges, every subsequence of it also converges to the same
limit.

The theorem~\ref{th:serie_positiva} tells us that series with positive terms are
regular and therefore satisfy the hypotheses of the theorem.
\end{proof}

\subsection{the harmonic series}

We observe that the ratio test does not apply to the \emph{harmonic series}%
\mymargin{harmonic series}%
\index{series!harmonic}%
\[
  \sum_k \frac{1}{k}
\]
since
\[
 \frac{\frac{1}{k+1}}{\frac{1}{k}}
 = \frac{k}{k+1} \to 1.
\]

To understand whether the harmonic series converges or diverges, we present the
method of \emph{condensation}, which will be stated in general in the next
theorem but can be better understood when applied to the particular case of the
harmonic series.

We will show that the harmonic series diverges.
The idea is simply to associate the terms of the harmonic series into groups of
length powers of two and estimate the sum of each group from below with the
smallest term (i.e., the last) of each group:
\begin{align*}
 \sum_{k=1}^{+\infty} \frac{1}{k}
 & = 1 + \frac 1 2
     + \enclose{\frac 1 3 + \frac 1 4}
     + \enclose{\frac 1 5 + \frac 1 6 + \frac 1 7 + \frac 1 8}
     + \dots\\
 & > 1 + \frac 1 2 + 2 \cdot \frac 1 4 + 4 \cdot \frac 1 8 + \dots \\
   & = 1 + \frac 1 2 + \frac 1 2 + \frac 1 2 + \dots
    = +\infty.
\end{align*}

\begin{theorem}[Cauchy's condensation test]%
\label{th:condensazione}%
\mymark{**}%
\mymargin{Cauchy's condensation test}%
\index{test!Cauchy's condensation}%
\index{condensation!Cauchy's test}%
\index{Cauchy!condensation test}%
Let $a_n$ be a decreasing sequence of non-negative real numbers: $a_n \ge 0$.
Then the series $\sum a_k$ converges if and only if the series
\[
  \sum 2^k a_{2^k}
\]
converges.
\end{theorem}
%
\begin{proof}
\mymark{**}
For convenience, suppose the sums start from $k=1$.
It involves grouping the terms $a_k$ into groups of powers of two:
\begin{align*}
  b_0 & = a_1, \\
  b_1 &= a_2 +  a_3, \\
  b_2 &= a_4 +  a_5 +  a_6 +  a_7, \\
  b_3 &= a_8 +  a_9 +  a_{10} + \dots + a_{15}, \\
  &\vdots\\
  b_n &= a_{2^n} +  a_{2^{n}+1} + \dots + a_{2^{n+1}-1},\\
  &\vdots
\end{align*}
Thanks to theorem~\ref{th:serie_associativa} on the associativity of series with
positive terms, we know that
  \[
  \sum_{k=1}^{+\infty} a_k = \sum_{n=0}^{+\infty} b_n.
  \]
Each term $b_n$ is the sum of $2^n$ terms of the sequence $a_k$, which, being a
decreasing sequence, can be estimated from above and below with the first and
last term of each sum, thus:
\[
  2^n a_{2^{n+1}-1} \le b_n \le 2^n a_{2^{n}}.
\]
By hypothesis, the series $\sum 2^n a_{2^n}$ is convergent, so by comparison,
the series $\sum b_n$ is also convergent.

Conversely, if the series $\sum 2^n a_{2^n}$ is divergent, then the series $\sum
2^{n+1} a_{2^{n+1}}$ is also divergent, and knowing that $b_n\ge 2^n
a_{2^{n+1}-1}\ge \frac 1 2 2^{n+1} a_{2^{n+1}}$, we obtain, by comparison, that
the series $\sum b_n$ is also divergent.
\end{proof}

\begin{corollary}[generalized harmonic series]
\mymark{***}
\mymargin{generalized harmonic series}%
\index{series!harmonic!generalized}
The series
\[
 \sum_n \frac{1}{n^\alpha}
\]
converges if $\alpha>1$,
diverges if $0\le \alpha\le 1$.
\end{corollary}
%
\begin{proof}
\mymark{***}
We apply the condensation test. Setting $a_n = \frac 1{n^\alpha}$, we have
\[
  \sum_n 2^n a_{2^n} = \sum_n 2^n \frac{1}{(2^n)^\alpha}
  = \sum_n 2^{n(1-\alpha)}
  = \sum_n \enclose{2^{(1-\alpha)}}^n
\]
which is a geometric series with ratio $q=2^{1-\alpha}$.
If $\alpha>1$, then $q<1$, and the harmonic series is convergent; if $\alpha \le
1$, then $q\ge 1$, and the harmonic series is divergent.
\end{proof}

\begin{exercise}
Use the condensation test to prove that the series
\[
  \sum \frac{1}{n \cdot \ln n}
\]
diverges.
\end{exercise}

\begin{exercise}
    Determine the values of the parameters $\alpha\in \RR$ and $\beta\in \RR$
  for which the following series is convergent:
  \[
    \sum \frac{1}{n^\alpha (\ln n)^\beta}.
  \]
\end{exercise}